\scanpage{30}
% archive foliation: 29
définition 1 on peut prendre le $K'$ une extension de \uline{type fini} de $K$.

Cela résulte de Chap V, § 1. Question : peut-on prendre pour $K'$ n'importe quelle extension \uline{\unsure{composée|commune}} ?

Cela montre que pour étudier les \deleted{\unsure{objets}} classes de faisceaux \inserted{sur $X$}, à équivalence près, on peut se borner aux faisceaux ``définis'' sur des extensions de type fini de $k$.
Cela prouve aussi \deleted{\unsure{que}} \inserted{qu'il est raisonnable de parler de classes} de classes de faisceaux à isom. près [qui est un ensemble dans le $\unsure{\mathcal U}$-univers \unsure{que} \unsure{alors|choisi} contenant $k,X$]. \inserted{$\unsure{\Phi|\mathrm{Cl}}(X/k)$}

\deleted{De plus, cet \unsure{ens} ne change pas par extension du corps de base $k$ en un $k'$. Enfin, $\unsure{\Phi|\mathrm{Cl}}(X/k)$ est un foncteur contravariant en $X\unsure{/k}$, pour $k$ fixé.}
\deleted{\uuline{Définition}}
\figure{\begin{tikzpicture}\draw (0,2)--(1.3,.1);\draw (2,2)..controls (2.5,1) and (3,.4)..(3.3,.2);\draw (4,2)..controls (4.4,1) and (4.8,.1)..(5,.2);\draw (5.8,2.3)--(6.4,.5);\draw (-.2,1.8)--(-.2,-.4);\draw (-.3,-.3)--(6.4,0);\end{tikzpicture}}
\margin{plus bas \unsure{après}}
On note que \unsure{l'ens} $\unsure{\Phi|\mathrm{Cl}}(X/k)$ est, pour \uline{$k$ fixé}, un foncteur contravariant en $X$. Cependant, si $k'$ est une extension de $k$, il n'y a pas en
